\documentclass{article}
\usepackage{amsmath,amssymb}
\begin{document}
Research notes for the matching-complex question.

Counterexample selected (rank 3 free group):
\[
w=\bar a a\, b\bar b\, \bar a a\, c\bar c\, \bar a a.
\]
It is reducible by adjacent cancellations and quasi-reduced, because the
separator pairs \(b\bar b\) and \(c\bar c\) break the forbidden
\(\ell\ell^{-1}\ell\) patterns that would occur in
\(\bar a a\bar a a\bar a a\).

Key structural facts used:
1. For any admissible matching, after inserting vertices at crossings,
the resulting planar graph is a forest.  Any graph cycle is interior
(boundary endpoints have valence one); an innermost cycle bounds a
region whose \(M\)-boundary component has no boundary endpoints, which
violates the definition.
2. If a generator occurs exactly once in each orientation, its two
endpoints are forced to be joined by an uncrossed interval.  At one
endpoint the two local sides must both end at the unique inverse-labelled
endpoint.  In a tree there is a unique path; at the first crossing on
that path the two sides diverge into different components, contradiction.
3. Therefore the \(b\bar b\) and \(c\bar c\) adjacent pairs in the
counterexample are fixed boundary-parallel intervals.  Removing them
identifies \(F_w\) with the analogous complex for
\(u=\bar a a\bar a a\bar a a\) on six endpoints.

Six-endpoint computation:
Odd endpoints are labelled \(\bar a\), even endpoints \(a\).  With three
intervals, admissible diagrams correspond exactly to endpoint pairings
whose interleaving graph is a forest.  The only non-forest pairing of
six points is \(14|25|36\), where all three pairs cross.  All other
fourteen pairings occur.  Parity lemma: for a face-side component from
\(i\) to \(j\), push it slightly into the face to an arc disjoint from
the diagram; each side of the cut contains an even number of endpoints,
so \(i,j\) have opposite parity, hence inverse labels for the alternating
word.

Cells for \(u\):
Vertices:
\(v_0=12|34|56\), \(v_1=12|36|45\), \(v_2=14|23|56\),
\(v_3=16|23|45\), \(v_4=16|25|34\).
Edges:
\(e_{01}=12|35|46\), \(e_{02}=13|24|56\), \(e_{13}=13|26|45\),
\(e_{23}=15|23|46\), \(e_{04}=15|26|34\), \(e_{34}=16|24|35\).
Squares:
\(Q_{12}=13|25|46\), \(Q_{14}=14|26|35\), \(Q_{24}=15|24|36\).
Boundaries:
\(Q_{12}\) has \(e_{01},e_{13},e_{23},e_{02}\);
\(Q_{14}\) has \(e_{01},e_{13},e_{34},e_{04}\);
\(Q_{24}\) has \(e_{02},e_{23},e_{34},e_{04}\).
The 1-skeleton is a theta graph with poles \(v_0,v_3\) and middle
vertices \(v_1,v_2,v_4\); the three squares fill the three theta cycles.
Hence the complex is \(S^2\), not contractible.

\end{document}
